\clearpage
\section{Prepare to explain and defend the work}\label{sec:explain}
\subsection{A short explanation you can adapt}
\begin{concept}{About one minute}
I am testing whether a particular way of training movement representations
makes matching human actions less sensitive to damaged pose information.
At evaluation, the system gets one clean example of each unfamiliar action
and must match other clips to those examples.

I reuse a pretrained MotionBERT encoder and keep it fixed. I train the same
linear comparison head with three recipes. The main comparison is the full
contextual recipe against standard contrastive training. Both use paired
starting weights, batches, training duration, and tuning opportunities.

I damage only the query poses, using fixed strengths of coordinate noise,
missing joints, and missing frames. The primary quantity is the difference
in clean-to-damaged accuracy drops, while absolute accuracy is reported too.

Historical pilots and input checks have passed. Current work is blocked on
storage and new-runtime verification. The final comparison has not run, so I
am not claiming that the contextual recipe wins.
\end{concept}

Say this in your own words. The structure matters more than memorizing it:
question, input/output, controlled comparison, measurement, current evidence.

\subsection{A five-minute technical explanation}
\begin{enumerate}
\item \textbf{Define the task.} A sequence of estimated 2D joint positions enters
a frozen encoder, then a learned head. The output is a unit vector used for
nearest-gallery retrieval. One-shot refers to one exemplar per novel action.
\item \textbf{Define the treatment.} Compare training recipes on the same head.
The contextual recipe includes neighborhoods, a changed margin, and a regularizer;
it is not an isolated neighborhood ablation.
\item \textbf{Explain the split and tuning.} Eighty actions train candidates;
twenty other auxiliary actions select rates using clean development retrieval.
Nine fresh final heads use all 100 auxiliary actions. Twenty novel actions
remain outside selection and training.
\item \textbf{Explain the intervention.} Three damage families at three fixed
strengths act after preprocessing and before encoding. Damage is paired across
methods and nested across severities; gallery examples remain clean.
\item \textbf{Explain the statistic.} Compare clean-to-corrupt Top-1 drops,
averaged over three fixed seeds and nine conditions. Resample whole performance
groups within actions to preserve camera-view dependence. State the interval's
conditional scope and show absolute accuracy beside the drop.
\item \textbf{Separate evidence from intent.} Describe completed historical
pilots and damage checks, then the current runtime/storage prerequisite. End
with what must happen before a final claim is possible.
\end{enumerate}

\subsection{Questions someone should be able to ask you}
\textbf{Why expect contextual learning to help?}
Relations to neighboring examples may provide useful information beyond a
pair's direct similarity. That is a plausible mechanism motivating the study.
Image-domain findings do not establish robustness to noisy pose queries; this
experiment tests that transfer rather than assuming it.

\textbf{What exactly is original here?}
The backbone, pose dataset, and published objective ideas are reused. The project
specifies and implements a controlled pose-retrieval comparison, its local
adaptations, corruption design, protected evaluation, and verification/reporting
process. It is not claiming to invent MotionBERT or contextual loss.

\textbf{Why is a linear head enough?}
The nonlinear encoder already supplies rich motion features. A common affine
head limits architectural variation while still learning a different geometry.
This is a deliberate comparison setting, not proof that a nonlinear head could
not perform better. Its approximately 17.8 million parameters also mean that
``linear'' should not be confused with negligible capacity.

\textbf{Does fixed training make the comparison perfectly fair?}
It makes the stated budget and controls comparable. It does not guarantee
equal computational cost, convergence, optimal hyperparameters, or equal
compatibility with every representation. The result is conditional on these
declared recipes and settings.

\textbf{Why not choose the best pilot or early checkpoint?}
Those outputs are engineering diagnostics. Reusing them for selection or choosing
an attractive epoch would give an extra, unplanned search opportunity and change
the fixed-duration procedure. The declared fresh selection runs have one rule.

\textbf{Why keep both official and stricter query sets?}
The official set enables comparison with the established benchmark convention.
The primary set removes other cameras' views of the gallery performance,
reducing a specific shortcut. Their distinct roles are fixed before results.

\textbf{Could the model learn who the person is?}
Action-disjoint splits do not remove all person or setup dependence. Shared
people, acquisition conditions, and representation biases remain limitations.
The study is not a separate test of generalization to unseen people.

\textbf{Why are three seeds not enough for a universal claim?}
They show three fixed training realizations and allow paired comparisons.
The primary bootstrap does not resample them and cannot infer performance over
all possible training randomness. Per-seed results show how consistent this
small set is without disguising its limited scope.

\textbf{What would count as success if contextual loses?}
A valid, complete, reproducible comparison with an honest contrary result.
Scientific success is answering the specified question. Failing to complete
the comparison because of capacity is a blocked study, not evidence of no effect.

\textbf{Why has verification taken so much work?}
The experiment needs more than a model that trains. Source/input identity,
exact pairing, damaged-input validity, runtime repeatability, and independently
reconstructible results determine whether the final comparison is interpretable.
The project also learned that those checks must be included in capacity planning.

\textbf{What cannot you claim yet?}
No final recipe ranking, novel-action robustness effect, or completed empirical
comparison. Successful pilots establish their execution; report drafts and a
working tracker do not supply missing final results.

\subsection{Draw these three things without looking}
\begin{enumerate}
\item The shape pipeline: $2\times100\times17\times3$ poses to frozen features,
8,704 pooled values, a learned 2,048-dimensional unit vector, and gallery ranking.
\item The action split: eighty candidate-training actions, twenty development
actions, all 100 used only after selection for fresh final training, and twenty
separate novel actions for final evaluation.
\item A two-row result table containing clean accuracy, corrupted accuracy,
and drop. Construct a case where contextual degrades less but has worse
corrupted accuracy, then explain why both facts belong in the report.
\end{enumerate}

\subsection{Use the notebooks to test understanding}
\begin{tabularx}{\linewidth}{@{}P{0.11\linewidth}Y Y@{}}
\toprule Lab & Concept & What you should explain afterward\\\midrule
01 & Cosine geometry and a normalized-head gradient & Why direction determines
ranking; how changing $W$ changes it; what finite differences verify.\\\addlinespace
02 & Splits, leakage, pairing, and equal selection budgets & Why repeated views
and unequal tuning opportunities can distort a comparison.\\\addlinespace
03 & Degradation and paired cluster bootstrap & Why smaller drop can coexist
with worse accuracy; why the sampled query count is the denominator.\\\bottomrule
\end{tabularx}

All examples are invented. Do not report their scores as project results.
Notebook 1 simplifies the training objective to make differentiation visible;
it does not implement or validate the contextual STE. Notebook 3 implements the
analysis logic on a tiny synthetic population, not the locked real-data campaign.

\appendix
\clearpage
\section{Answers to the study checks}
\begin{enumerate}[label=\textbf{\arabic*.},leftmargin=2em]
\item \textbf{Random row splits.} They can put the same action classes in
training and evaluation, testing new clips of familiar actions rather than
transfer to unfamiliar actions. They can also separate synchronized views of
one performance across roles.
\item \textbf{Gallery exclusion.} A different camera file can still depict the
same event as the gallery example. The primary set excludes the underlying
performance group, not only one filename.
\item \textbf{Learning with a frozen encoder.} The encoder's output stays fixed,
but the head's $W,b$ change how those features map into comparison directions.
A novel action does not need a learned output-class coordinate: it supplies a
gallery exemplar and is recognized by similarity to it.
\item \textbf{Cache perturbation.} A pose corruption is transformed nonlinearly
by the encoder. In general $F(T(X))$ is not equivalent to adding noise to $F(X)$.
The intervention location is therefore part of the hypothesis.
\item \textbf{Physical batch.} Pairwise similarities, contextual neighbors, and
SupCon denominators depend on the examples simultaneously present. Summing
gradients from separate smaller batches omits cross-microbatch relations.
\item \textbf{Full recipe.} Contextual and standard contrastive differ in the
neighborhood term, positive margin, and similarity regularizer. A matched
ablation would be needed to isolate one component's causal contribution.
\item \textbf{Loss magnitudes.} The formulas, reductions, coefficients, and
attainable minima differ. Compare retrieval under the common evaluation,
not which method reports the smallest numerical training loss.
\item \textbf{Shared severity randomness.} Reuse of the underlying noise field
or nested mask makes a stronger level a controlled strengthening of the same
damage. Unrelated draws could obscure the effect of increasing severity.
\item \textbf{Fallback and realism.} A training-only fallback avoids estimating
preprocessing parameters from held-out poses. Synthetic zero masks and Gaussian
noise do not reproduce every output a real pose estimator produces under camera
occlusion, so they support a narrower robustness claim.
\item \textbf{Groups versus query weights.} Groups preserve dependence during
resampling; query weights preserve the chosen statistic. Equal weighting of
unequal-sized group means would change the quantity being estimated.
\item \textbf{Interval crossing zero.} The data do not support the prespecified
smaller-drop claim at that rule. This is neither proof of equality nor an
estimate that the probability of no effect is 95\%.
\item \textbf{Bootstrap draws and seeds.} Repeatedly resampling performance
groups recombines evidence from the same three heads per method. It creates no
new independent training realizations or uncertainty over unobserved seeds.
\end{enumerate}

\section{Notation and essential terms}
\begin{longtable}{@{}P{0.25\linewidth}P{0.68\linewidth}@{}}
\toprule Symbol or term & Meaning here\\\midrule\endhead
$X$ & Prepared pose tensor for one capture.\\
$h\in\R^{8704}$ & Frozen, pooled, unnormalized pre-head feature vector.\\
$W,b$ & Trainable affine head weights and bias.\\
$u\in\R^{2048}$; $z$ & Head output before normalization; its unit-length direction.\\
$s_{ij}$; $d_{ij}$ & Cosine similarity; squared unit-vector distance.\\
$y_{ij}$ & Indicator that two training examples have the same action label.\\
$P,K,n$ & Classes per batch, examples per class, and physical batch size.\\
$k,\epsilon,\alpha$ & Contextual neighbor rank, membership tolerance, and
surrogate-gradient factor.\\
$\tau$; $\eta$ & SupCon temperature; optimizer learning rate.\\
$A_{m,s,c}$; $d_{m,s,c}$ & Top-1 accuracy and clean-to-corrupt drop; this $d$ has
method/seed/condition indices, unlike pairwise distance $d_{ij}$.\\
$\Delta$ & Contextual drop minus contrastive drop, averaged over fixed seeds
and corruptions. Negative favors less contextual degradation.\\
Gallery; query & Labeled reference examples; the example to be matched.\\
Embedding & A learned numerical representation used for comparison.\\
Epoch & The prescribed number of balanced updates, not necessarily one visit
to each row.\\
Hyperparameter & A setting chosen outside gradient updates, such as learning
rate; distinct from learned weights.\\
Leakage & Information from an evaluation role influences development or makes
the evaluation easier in a way the claimed design excludes.\\
Estimand & The exact quantity the experiment intends to estimate.\\
Bootstrap & Recompute a statistic over repeated resamples of the declared units.\\
Provenance & The source, input, configuration, and runtime identities connecting
an output to the execution that produced it.\\
Gate & An explicit requirement and its acceptance evidence before later work
is eligible.\\\bottomrule
\end{longtable}

\section{Source map and evidence boundaries}\label{sec:sources}
The linked repository sources are pinned to commit
\code{b32d50f66206b561fa34867c9274fac6ba7c287d}, the main-branch record through
October 25, unless marked as the October 28 preparation record at
\code{5d50410b07e19a980fe3b711e63b10417b40b818}. This keeps the account traceable
even if a live branch later changes. The guide is explanatory documentation;
the canonical protocol and authorized amendments control execution.

\begin{longtable}{@{}P{0.29\linewidth}P{0.65\linewidth}@{}}
\toprule Read this for & Primary project source\\\midrule\endhead
Scientific question and rules & \src{plan/research-plan.v3.md}{Final v3 plan}
and \src{docs/protocol/protocol-v3.md}{protocol v3}.\\\addlinespace
Machine-readable settings & \src{configs/protocol.v3.yaml}{Protocol template}
and \href{https://github.com/JJCAPPE/pose-embedding/tree/b32d50f66206b561fa34867c9274fac6ba7c287d/configs/experiments/v3}{27 candidate configurations}. A template is
not an executed or resolved design lock.\\\addlinespace
Method equations and adaptation & \src{docs/protocol/contextual-v3-source-notes.md}{Contextual source notes},
\src{src/pose_embed/losses/contextual.py}{contextual loss}, and
\src{src/pose_embed/losses/pairwise.py}{pairwise losses}.\\\addlinespace
Reuse and preprocessing & \src{third_party/upstreams.toml}{Upstream manifest},
\src{THIRD_PARTY_NOTICES.md}{notices}, and
\src{docs/protocol/motionbert-adapter-gate.md}{adapter acceptance}.\\\addlinespace
Why the scope changed & \src{docs/decisions/0004-frozen-one-shot-v3.md}{Decision 0004};
\src{docs/decisions/0006-v3-execution-completion-amendment.md}{0006},
\src{docs/decisions/0007-safe-validation-reuse-and-stage-capacity.md}{0007},
\src{docs/decisions/0008-runtime-baseline.md}{0008}, and
\src{docs/decisions/0009-updated-kernel-baseline.md}{0009}.\\\addlinespace
Accepted history and current state & \src{docs/protocol/week-4-execution.md}{Week 4},
\src{docs/protocol/week-5-execution.md}{Week 5 through October 25}, and
\prepsrc{docs/protocol/week-5-execution.md}{October 28 current checkpoint}.\\\bottomrule
\end{longtable}

The October 28 schedule's opening prose overstates the October 25 filesystem
failure. The numerical Week 5 record shows quota below the requirement while
filesystem free space exceeded it; usable space was still insufficient because
the minimum governs. This guide uses the numerical record and states the
storage blocker without treating an old observation as live capacity.

The intended prerequisites are general calculus, linear algebra, elementary
probability/statistics, and strong programming. No transcript, grades, student
identifier, or private academic record is included. Figures are schematic and
notebook data are synthetic. No licensed pose array, learned checkpoint,
participant record, or private operational log is included.

\begin{thebibliography}{9}
\addcontentsline{toc}{section}{Research origins}
\bibitem{ntu} J. Liu et al. \emph{NTU RGB+D 120: A Large-Scale Benchmark for
3D Human Activity Understanding}. IEEE TPAMI, 2020.
\href{https://arxiv.org/abs/1905.04757}{Authors' paper};
\href{https://rose1.ntu.edu.sg/dataset/actionRecognition/}{Official dataset}.
The local input contract separately identifies the adopted HRNet-derived 2D
annotations; the dataset's title does not make our input coordinates 3D.
\bibitem{motionbert} W. Zhu et al. \emph{MotionBERT: A Unified Perspective on
Learning Human Motion Representations}. ICCV, 2023.
\href{https://openaccess.thecvf.com/content/ICCV2023/html/Zhu_MotionBERT_A_Unified_Perspective_on_Learning_Human_Motion_Representations_ICCV_2023_paper.html}{Official paper and supplement}.
Read for the encoder's original purpose; our frozen-head comparison is a
different experiment from its downstream fine-tuning results.
\bibitem{contextual} C. Liao, T. Tsiligkaridis, and B. Kulis.
\emph{Supervised Metric Learning to Rank for Retrieval via Contextual Similarity
Optimization}. ICML, 2023.
\href{https://proceedings.mlr.press/v202/liao23b.html}{Paper and appendix}.
Read for the contextual objective's origin. The local source notes resolve
equation/backward conventions and document pose-study adaptations.
\bibitem{supcon} P. Khosla et al. \emph{Supervised Contrastive Learning}.
NeurIPS, 2020.
\href{https://papers.nips.cc/paper/2020/hash/d89a66c7c80a29b1bdbab0f2a1a94af8-Abstract.html}{Official paper}.
Read for label-based contrastive learning; this study fixes a single-view variant.
\end{thebibliography}
